% Options for packages loaded elsewhere
% Options for packages loaded elsewhere
\PassOptionsToPackage{unicode}{hyperref}
\PassOptionsToPackage{hyphens}{url}
%
\documentclass[
  english,
  russian,
  12pt,
  a4paper,
  DIV=11,
  numbers=noendperiod]{scrreprt}
\usepackage{xcolor}
\usepackage{amsmath,amssymb}
\setcounter{secnumdepth}{5}
\usepackage{iftex}
\ifPDFTeX
  \usepackage[T1]{fontenc}
  \usepackage[utf8]{inputenc}
  \usepackage{textcomp} % provide euro and other symbols
\else % if luatex or xetex
  \usepackage{unicode-math} % this also loads fontspec
  \defaultfontfeatures{Scale=MatchLowercase}
  \defaultfontfeatures[\rmfamily]{Ligatures=TeX,Scale=1}
\fi
\usepackage{lmodern}
\ifPDFTeX\else
  % xetex/luatex font selection
\fi
% Use upquote if available, for straight quotes in verbatim environments
\IfFileExists{upquote.sty}{\usepackage{upquote}}{}
\IfFileExists{microtype.sty}{% use microtype if available
  \usepackage[]{microtype}
  \UseMicrotypeSet[protrusion]{basicmath} % disable protrusion for tt fonts
}{}
\usepackage{setspace}
% Make \paragraph and \subparagraph free-standing
\makeatletter
\ifx\paragraph\undefined\else
  \let\oldparagraph\paragraph
  \renewcommand{\paragraph}{
    \@ifstar
      \xxxParagraphStar
      \xxxParagraphNoStar
  }
  \newcommand{\xxxParagraphStar}[1]{\oldparagraph*{#1}\mbox{}}
  \newcommand{\xxxParagraphNoStar}[1]{\oldparagraph{#1}\mbox{}}
\fi
\ifx\subparagraph\undefined\else
  \let\oldsubparagraph\subparagraph
  \renewcommand{\subparagraph}{
    \@ifstar
      \xxxSubParagraphStar
      \xxxSubParagraphNoStar
  }
  \newcommand{\xxxSubParagraphStar}[1]{\oldsubparagraph*{#1}\mbox{}}
  \newcommand{\xxxSubParagraphNoStar}[1]{\oldsubparagraph{#1}\mbox{}}
\fi
\makeatother


\usepackage{longtable,booktabs,array}
\usepackage{calc} % for calculating minipage widths
% Correct order of tables after \paragraph or \subparagraph
\usepackage{etoolbox}
\makeatletter
\patchcmd\longtable{\par}{\if@noskipsec\mbox{}\fi\par}{}{}
\makeatother
% Allow footnotes in longtable head/foot
\IfFileExists{footnotehyper.sty}{\usepackage{footnotehyper}}{\usepackage{footnote}}
\makesavenoteenv{longtable}
\usepackage{graphicx}
\makeatletter
\newsavebox\pandoc@box
\newcommand*\pandocbounded[1]{% scales image to fit in text height/width
  \sbox\pandoc@box{#1}%
  \Gscale@div\@tempa{\textheight}{\dimexpr\ht\pandoc@box+\dp\pandoc@box\relax}%
  \Gscale@div\@tempb{\linewidth}{\wd\pandoc@box}%
  \ifdim\@tempb\p@<\@tempa\p@\let\@tempa\@tempb\fi% select the smaller of both
  \ifdim\@tempa\p@<\p@\scalebox{\@tempa}{\usebox\pandoc@box}%
  \else\usebox{\pandoc@box}%
  \fi%
}
% Set default figure placement to htbp
\def\fps@figure{htbp}
\makeatother



\ifLuaTeX
\usepackage[bidi=basic,provide=*]{babel}
\else
\usepackage[bidi=default,provide=*]{babel}
\fi
% get rid of language-specific shorthands (see #6817):
\let\LanguageShortHands\languageshorthands
\def\languageshorthands#1{}


\setlength{\emergencystretch}{3em} % prevent overfull lines

\providecommand{\tightlist}{%
  \setlength{\itemsep}{0pt}\setlength{\parskip}{0pt}}



 
\usepackage[style=gost-numeric,backend=biber,langhook=extras,autolang=other*]{biblatex}
\addbibresource{bib/cite.bib}

\usepackage[]{csquotes}

\usepackage{indentfirst}
\usepackage{float}
\floatplacement{figure}{H}
\IfFileExists{plex-otf.sty}{
  %% Full TeXlive
  \usepackage[
    % math,
    RM={Scale=0.94},SS={Scale=0.94},SScon={Scale=0.94},TT={Scale=MatchLowercase,FakeStretch=0.9},
    DefaultFeatures={Ligatures=Common}
  ]{plex-otf}
  % \usepackage[Scale=MatchUppercase]{juliamono}
}{
  %% TinyTeX
  \usepackage{libertine}
}
\KOMAoption{captions}{tableheading}
\makeatletter
\@ifpackageloaded{caption}{}{\usepackage{caption}}
\AtBeginDocument{%
\ifdefined\contentsname
  \renewcommand*\contentsname{Содержание}
\else
  \newcommand\contentsname{Содержание}
\fi
\ifdefined\listfigurename
  \renewcommand*\listfigurename{Список иллюстраций}
\else
  \newcommand\listfigurename{Список иллюстраций}
\fi
\ifdefined\listtablename
  \renewcommand*\listtablename{Список таблиц}
\else
  \newcommand\listtablename{Список таблиц}
\fi
\ifdefined\figurename
  \renewcommand*\figurename{Рисунок}
\else
  \newcommand\figurename{Рисунок}
\fi
\ifdefined\tablename
  \renewcommand*\tablename{Таблица}
\else
  \newcommand\tablename{Таблица}
\fi
}
\@ifpackageloaded{float}{}{\usepackage{float}}
\floatstyle{ruled}
\@ifundefined{c@chapter}{\newfloat{codelisting}{h}{lop}}{\newfloat{codelisting}{h}{lop}[chapter]}
\floatname{codelisting}{Список}
\newcommand*\listoflistings{\listof{codelisting}{Листинги}}
\makeatother
\makeatletter
\makeatother
\makeatletter
\@ifpackageloaded{caption}{}{\usepackage{caption}}
\@ifpackageloaded{subcaption}{}{\usepackage{subcaption}}
\makeatother
\usepackage{bookmark}
\IfFileExists{xurl.sty}{\usepackage{xurl}}{} % add URL line breaks if available
\urlstyle{same}
\hypersetup{
  pdftitle={Первоначальна настройка git},
  pdfauthor={Смирнов Артём Сергеевич},
  pdflang={ru-RU},
  hidelinks,
  pdfcreator={LaTeX via pandoc}}


\title{Первоначальна настройка git}
\usepackage{etoolbox}
\makeatletter
\providecommand{\subtitle}[1]{% add subtitle to \maketitle
  \apptocmd{\@title}{\par {\large #1 \par}}{}{}
}
\makeatother
\subtitle{Архитектура Компьютеров}
\author{Смирнов Артём Сергеевич}
\date{}
\begin{document}
\maketitle

\renewcommand*\contentsname{Содержание}
{
\setcounter{tocdepth}{1}
\tableofcontents
}
\listoffigures
\listoftables

\setstretch{1.5}
\chapter{Цель
работы}\label{ux446ux435ux43bux44c-ux440ux430ux431ux43eux442ux44b}

\begin{enumerate}
\def\labelenumi{\arabic{enumi})}
\tightlist
\item
  Изучить идеологию и применение средств контроля версий.
\item
  Освоить умения по работе с git.
\end{enumerate}

\chapter{Задание}\label{ux437ux430ux434ux430ux43dux438ux435}

\begin{enumerate}
\def\labelenumi{\arabic{enumi})}
\tightlist
\item
  Создать базовую конфигурацию для работы с git.
\item
  Создать ключ SSH.
\item
  Создать ключ PGP.
\item
  Настроить подписи git.
\item
  Зарегистрироваться на Github.
\item
  Создать локальный каталог для выполнения заданий по предмету.
\end{enumerate}

\chapter{Теоретическое
введение}\label{ux442ux435ux43eux440ux435ux442ux438ux447ux435ux441ux43aux43eux435-ux432ux432ux435ux434ux435ux43dux438ux435}

\section{Основные команды
git}\label{ux43eux441ux43dux43eux432ux43dux44bux435-ux43aux43eux43cux430ux43dux434ux44b-git}

\subsection{Перечислим наиболее часто используемые команды
git.}\label{ux43fux435ux440ux435ux447ux438ux441ux43bux438ux43c-ux43dux430ux438ux431ux43eux43bux435ux435-ux447ux430ux441ux442ux43e-ux438ux441ux43fux43eux43bux44cux437ux443ux435ux43cux44bux435-ux43aux43eux43cux430ux43dux434ux44b-git.}

\subsubsection{Создание основного дерева
репозитория:}\label{ux441ux43eux437ux434ux430ux43dux438ux435-ux43eux441ux43dux43eux432ux43dux43eux433ux43e-ux434ux435ux440ux435ux432ux430-ux440ux435ux43fux43eux437ux438ux442ux43eux440ux438ux44f}

git init Получение обновлений (изменений) текущего дерева из
центрального репозитория:

git pull Отправка всех произведённых изменений локального дерева в
центральный репозиторий:

git push Просмотр списка изменённых файлов в текущей директории:

git status Просмотр текущих изменений:

git diff Сохранение текущих изменений:

добавить все изменённые и/или созданные файлы и/или каталоги:

git add . добавить конкретные изменённые и/или созданные файлы и/или
каталоги:

git add имена\_файлов удалить файл и/или каталог из индекса репозитория
(при этом файл и/или каталог остаётся в локальной директории):

git rm имена\_файлов Сохранение добавленных изменений:

сохранить все добавленные изменения и все изменённые файлы:

git commit -am \enquote*{Описание коммита} сохранить добавленные
изменения с внесением комментария через встроенный редактор:

git commit создание новой ветки, базирующейся на текущей:

git checkout -b имя\_ветки переключение на некоторую ветку:

git checkout имя\_ветки (при переключении на ветку, которой ещё нет в
локальном репозитории, она будет создана и связана с удалённой) отправка
изменений конкретной ветки в центральный репозиторий:

git push origin имя\_ветки слияние ветки с текущим деревом:

git merge --no-ff имя\_ветки Удаление ветки:

удаление локальной уже слитой с основным деревом ветки:

git branch -d имя\_ветки принудительное удаление локальной ветки:

git branch -D имя\_ветки удаление ветки с центрального репозитория:

git push origin :имя\_ветки

\subsubsection{Стандартные процедуры работы при наличии центрального
репозитория}\label{ux441ux442ux430ux43dux434ux430ux440ux442ux43dux44bux435-ux43fux440ux43eux446ux435ux434ux443ux440ux44b-ux440ux430ux431ux43eux442ux44b-ux43fux440ux438-ux43dux430ux43bux438ux447ux438ux438-ux446ux435ux43dux442ux440ux430ux43bux44cux43dux43eux433ux43e-ux440ux435ux43fux43eux437ux438ux442ux43eux440ux438ux44f}

Работа пользователя со своей веткой начинается с проверки и получения
изменений из центрального репозитория (при этом в локальное дерево до
начала этой процедуры не должно было вноситься изменений):

git checkout master git pull git checkout -b имя\_ветки Затем можно
вносить изменения в локальном дереве и/или ветке.

После завершения внесения какого-то изменения в файлы и/или каталоги
проекта необходимо разместить их в центральном репозитории. Для этого
необходимо проверить, какие файлы изменились к текущему моменту:

git status При необходимости удаляем лишние файлы, которые не хотим
отправлять в центральный репозиторий.

Затем полезно просмотреть текст изменений на предмет соответствия
правилам ведения чистых коммитов:

git diff Если какие-либо файлы не должны попасть в коммит, то помечаем
только те файлы, изменения которых нужно сохранить. Для этого используем
команды добавления и/или удаления с нужными опциями:

git add \ldots{}\\
git rm \ldots{} Если нужно сохранить все изменения в текущем каталоге,
то используем:

git add . Затем сохраняем изменения, поясняя, что было сделано:

git commit -am \enquote{Some commit message} Отправляем изменения в
центральный репозиторий:

git push origin имя\_ветки или

git push

Работа с локальным репозиторием Создадим локальный репозиторий. Сначала
сделаем предварительную конфигурацию, указав имя и email владельца
репозитория:

git config --global user.name \enquote{Имя Фамилия} git config --global
user.email \enquote{work@mail} Настроим utf-8 в выводе сообщений git:

git config --global quotepath false Для инициализации локального
репозитория, расположенного, например, в каталоге
\textasciitilde/tutorial, необходимо ввести в командной строке:

cd mkdir tutorial cd tutorial git init После это в каталоге tutorial
появится каталог .git, в котором будет храниться история изменений.
Создадим тестовый текстовый файл hello.txt и добавим его в локальный
репозиторий:

echo \enquote*{hello world} \textgreater{} hello.txt git add hello.txt
git commit -am \enquote*{Новый файл} Воспользуемся командой status для
просмотра изменений в рабочем каталоге, сделанных с момента последней
ревизии:

git status Во время работы над проектом так или иначе могут создаваться
файлы, которые не требуется добавлять в последствии в репозиторий.
Например, временные файлы, создаваемые редакторами, или объектные файлы,
создаваемые компиляторами. Можно прописать шаблоны игнорируемых при
добавлении в репозиторий типов файлов в файл .gitignore с помощью
сервисов. Для этого сначала нужно получить список имеющихся шаблонов:

curl -L -s https://www.gitignore.io/api/list Затем скачать шаблон,
например, для C и C++

curl -L -s https://www.gitignore.io/api/c \textgreater\textgreater{}
.gitignore curl -L -s https://www.gitignore.io/api/c++
\textgreater\textgreater{} .gitignore

\subsubsection{Работа с сервером
репозиториев}\label{ux440ux430ux431ux43eux442ux430-ux441-ux441ux435ux440ux432ux435ux440ux43eux43c-ux440ux435ux43fux43eux437ux438ux442ux43eux440ux438ux435ux432}

Для последующей идентификации пользователя на сервере репозиториев
необходимо сгенерировать пару ключей (приватный и открытый):

ssh-keygen -C \enquote{Имя Фамилия
\href{mailto:work@mail}{\nolinkurl{work@mail}}} Ключи сохраняться в
каталоге \textasciitilde/.ssh/. Существует несколько доступных серверов
репозиториев с возможностью бесплатного размещения данных. Например,
https://github.com/. Для работы с ним необходимо сначала завести на
сайте https://github.com/ учётную запись. Затем необходимо загрузить
сгенерённый нами ранее открытый ключ. Для этого зайти на сайт
https://github.com/ под своей учётной записью и перейти в меню GitHub
setting. После этого выбрать в боковом меню GitHub
setting\textgreater SSH-ключи и нажать кнопку Добавить ключ. Скопировав
из локальной консоли ключ в буфер обмена:

cat \textasciitilde/.ssh/id\_rsa.pub \textbar{} xclip -sel clip
Вставляем ключ в появившееся на сайте поле.

После этого можно создать на сайте репозиторий, выбрав в меню , дать ему
название и сделать общедоступным (публичным).

Для загрузки репозитория из локального каталога на сервер выполняем
следующие команды:

git remote add origin ssh://git@github.com//.git git push -u origin
master Далее на локальном компьютере можно выполнять стандартные
процедуры для работы с git при наличии центрального репозитория.

\subsection{Базовая настройка
git}\label{ux431ux430ux437ux43eux432ux430ux44f-ux43dux430ux441ux442ux440ux43eux439ux43aux430-git}

\subsubsection{Первичная настройка параметров
git}\label{ux43fux435ux440ux432ux438ux447ux43dux430ux44f-ux43dux430ux441ux442ux440ux43eux439ux43aux430-ux43fux430ux440ux430ux43cux435ux442ux440ux43eux432-git}

Зададим имя и email владельца репозитория:

git config --global user.name \enquote{Name Surname} git config --global
user.email \enquote{work@mail} Настроим utf-8 в выводе сообщений git:

git config --global core.quotepath false Настройте верификацию и
подписание коммитов git. Зададим имя начальной ветки (будем называть её
master):

git config --global init.defaultBranch master

git ls-files --eol Параметр autocrlf

Настройка core.autocrlf предназначена для того, чтобы в главном
репозитории все переводы строк текстовых файлах были одинаковы.
Настройка core.autocrlf с параметрами true и input делает все переводы
строк текстовых файлов в главном репозитории одинаковыми. core.autocrlf
true: конвертация CRLF-\textgreater LF при коммите и обратно
LF-\textgreater CRLF при выгрузке кода из репозитория на файловую
систему (обычно используется в Windows). core.autocrlf input:
конвертация CRLF-\textgreater LF только при коммитах (используются в
MacOS/Linux).

\subsection{Параметр
safecrlf}\label{ux43fux430ux440ux430ux43cux435ux442ux440-safecrlf}

Настройка core.safecrlf предназначена для проверки, является ли
окончаний строк обратимым для текущей настройки core.autocrlf.
core.safecrlf true: запрещается необратимое преобразование
lf\textless-\textgreater crlf. Полезно, когда существуют бинарные файлы,
похожие на текстовые файлы. core.safecrlf warn: печать предупреждения,
но коммиты с необратимым переходом принимаются. Установка параметра:

git config --global core.safecrlf warn

\section{Создание ключа
ssh}\label{ux441ux43eux437ux434ux430ux43dux438ux435-ux43aux43bux44eux447ux430-ssh}

Ключ ssh создаётся командой:

ssh-keygen -t Создайте ключи:

по алгоритму rsa с ключём размером 4096 бит:

ssh-keygen -t rsa -b 4096 по алгоритму ed25519:

ssh-keygen -t ed25519 При создании ключа команда попросит ввести любую
ключевую фразу для более надёжной защиты вашего пароля. Можно пропустить
этот этап, нажав Enter. Сменить пароль на ключ можно с помощью команды:

ssh-keygen -p

\section{Добавление SSH-ключа в учётную запись
GitHub}\label{ux434ux43eux431ux430ux432ux43bux435ux43dux438ux435-ssh-ux43aux43bux44eux447ux430-ux432-ux443ux447ux451ux442ux43dux443ux44e-ux437ux430ux43fux438ux441ux44c-github}

Скопируйте созданный SSH-ключ в буфер обмена командой:

xclip -i \textless{} \textasciitilde/.ssh/id\_ed25519.pub Откройте
настройки своего аккаунта на GitHub и перейдем в раздел SSH and GPC
keys. Нажмите кнопку ew SSH key. Добавьте в поле Title название этого
ключа, например, ed25519@hostname. Вставьте из буфера обмена в поле Key
ключ. Нажмите кнопку Add SSH key.

\chapter{Выполнение лабораторной
работы}\label{ux432ux44bux43fux43eux43bux43dux435ux43dux438ux435-ux43bux430ux431ux43eux440ux430ux442ux43eux440ux43dux43eux439-ux440ux430ux431ux43eux442ux44b}

\section{Порядок выполнения лабораторной
работы}\label{ux43fux43eux440ux44fux434ux43eux43a-ux432ux44bux43fux43eux43bux43dux435ux43dux438ux44f-ux43bux430ux431ux43eux440ux430ux442ux43eux440ux43dux43eux439-ux440ux430ux431ux43eux442ux44b}

\begin{figure}

\centering{

\includegraphics[width=0.7\linewidth,height=\textheight,keepaspectratio]{image/1.png}

}

\caption{\label{fig-001}Установка git}

\end{figure}%

\begin{figure}

\centering{

\includegraphics[width=0.7\linewidth,height=\textheight,keepaspectratio]{image/2.png}

}

\caption{\label{fig-002}Установка gh}

\end{figure}%

\begin{figure}

\centering{

\includegraphics[width=0.7\linewidth,height=\textheight,keepaspectratio]{image/3.png}

}

\caption{\label{fig-003}Введение имя пользователя и почты владельца
репозитория}

\end{figure}%

\begin{figure}

\centering{

\includegraphics[width=0.7\linewidth,height=\textheight,keepaspectratio]{image/4.png}

}

\caption{\label{fig-004}Настройка вывода сообщения формата utf-8}

\end{figure}%

\begin{figure}

\centering{

\includegraphics[width=0.7\linewidth,height=\textheight,keepaspectratio]{image/6.png}

}

\caption{\label{fig-005}Наименование ветки именем master}

\end{figure}%

\begin{figure}

\centering{

\includegraphics[width=0.7\linewidth,height=\textheight,keepaspectratio]{image/7.png}

}

\caption{\label{fig-006}Настраиваем параметры autocrlf и safecrlf}

\end{figure}%

\begin{figure}

\centering{

\includegraphics[width=0.7\linewidth,height=\textheight,keepaspectratio]{image/8.png}

}

\caption{\label{fig-007}Создание SSH ключа размерностью 4096 бит}

\end{figure}%

\begin{figure}

\centering{

\includegraphics[width=0.7\linewidth,height=\textheight,keepaspectratio]{image/9.png}

}

\caption{\label{fig-008}Создание SSH ключа}

\end{figure}%

\begin{figure}

\centering{

\includegraphics[width=0.7\linewidth,height=\textheight,keepaspectratio]{image/10.png}

}

\caption{\label{fig-009}Создание PGP ключа}

\end{figure}%

\begin{figure}

\centering{


\includegraphics[width=0.7\linewidth,height=\textheight,keepaspectratio]{image/11.png}

}

\caption{\label{fig-010}Настройка PGP ключа (Выбор типа ключа,
размерности ключа и срок его действия)}

\end{figure}%

\begin{figure}

\centering{


\includegraphics[width=0.7\linewidth,height=\textheight,keepaspectratio]{image/12.png}

}

\caption{\label{fig-001}Ввод личных данных(Имя и почта)}

\end{figure}%

\begin{figure}

\centering{

\includegraphics[width=0.7\linewidth,height=\textheight,keepaspectratio]{image/13.png}

}

\caption{\label{fig-0011}Соглашение на корректность заполненной
информации о ключе PGP}

\end{figure}%

\begin{figure}

\centering{

\includegraphics[width=0.7\linewidth,height=\textheight,keepaspectratio]{image/14.png}

}

\caption{\label{fig-0012}Вводим пароль, для ввода личных данных}

\end{figure}%

\begin{figure}

\centering{

\includegraphics[width=0.7\linewidth,height=\textheight,keepaspectratio]{image/git.png}

}

\caption{\label{fig-0013}Добавление личной информации в профиль на
GitHub}

\end{figure}%

\begin{figure}

\centering{

\includegraphics[width=0.7\linewidth,height=\textheight,keepaspectratio]{image/15.png}

}

\caption{\label{fig-0014}Узнаем отпечаток PGP ключа и копируем ключ PGP
в буфер обмен}

\end{figure}%

\begin{figure}

\centering{

\includegraphics[width=0.7\linewidth,height=\textheight,keepaspectratio]{image/16.png}

}

\caption{\label{fig-0015}Настройка автоматических подписей коммитов git}

\end{figure}%

\begin{figure}

\centering{

\includegraphics[width=0.7\linewidth,height=\textheight,keepaspectratio]{image/17.png}

}

\caption{\label{fig-0016}Настройка gh}

\end{figure}%

\begin{figure}

\centering{

\includegraphics[width=0.7\linewidth,height=\textheight,keepaspectratio]{image/18.png}

}

\caption{\label{fig-0017}Создаем директорию \enquote{Операционные
системы} и переход в этот каталог}

\end{figure}%

\begin{figure}

\centering{

\includegraphics[width=0.7\linewidth,height=\textheight,keepaspectratio]{image/19.png}

}

\caption{\label{fig-0018}Выполняем создание копии репозитория на основе
шаблона}

\end{figure}%

\begin{figure}

\centering{

\includegraphics[width=0.7\linewidth,height=\textheight,keepaspectratio]{image/20.png}

}

\caption{\label{fig-0019}Клонируем репозитории в созданную директорию в
папку os-intro}

\end{figure}%

\begin{figure}

\centering{

\includegraphics[width=0.7\linewidth,height=\textheight,keepaspectratio]{image/21.png}

}

\caption{\label{fig-00120}Переход в каталог os-intro, удаляем лишний
файл и создаем файл COURSE}

\end{figure}%

\begin{figure}

\centering{

\includegraphics[width=0.7\linewidth,height=\textheight,keepaspectratio]{image/22.png}

}

\caption{\label{fig-00121}Пишем команду make для создания новых, нужных
для прохождения курса, директорий}

\end{figure}%

\begin{figure}

\centering{

\includegraphics[width=0.7\linewidth,height=\textheight,keepaspectratio]{image/23.png}

}

\caption{\label{fig-00221}Мы обнародываем новые файлы для передачи}

\end{figure}%

\begin{figure}

\centering{

\includegraphics[width=0.7\linewidth,height=\textheight,keepaspectratio]{image/24.png}

}

\caption{\label{fig-00231}Передача новых файлов и каталогов на GitHub}

\end{figure}%

\chapter{Ответ на контрольные
вопросы}\label{ux43eux442ux432ux435ux442-ux43dux430-ux43aux43eux43dux442ux440ux43eux43bux44cux43dux44bux435-ux432ux43eux43fux440ux43eux441ux44b}

\section{Что такое системы контроля версий (VCS) и для решения каких
задач они
предназначаются?}\label{ux447ux442ux43e-ux442ux430ux43aux43eux435-ux441ux438ux441ux442ux435ux43cux44b-ux43aux43eux43dux442ux440ux43eux43bux44f-ux432ux435ux440ux441ux438ux439-vcs-ux438-ux434ux43bux44f-ux440ux435ux448ux435ux43dux438ux44f-ux43aux430ux43aux438ux445-ux437ux430ux434ux430ux447-ux43eux43dux438-ux43fux440ux435ux434ux43dux430ux437ux43dux430ux447ux430ux44eux442ux441ux44f}

Система контроля версий (Version Control System, VCS) --- это
программный инструмент, помогающий разработчикам (и не только)
отслеживать и управлять изменениями в коде (или любых других файлах) с
течением времени.

Основные задачи, которые решает VCS:

\subsection{Отслеживание истории изменений: Сохраняется информация о
том, кто, когда и какие именно изменения внес в файлы. При необходимости
можно посмотреть, как файл выглядел в
прошлом.}\label{ux43eux442ux441ux43bux435ux436ux438ux432ux430ux43dux438ux435-ux438ux441ux442ux43eux440ux438ux438-ux438ux437ux43cux435ux43dux435ux43dux438ux439-ux441ux43eux445ux440ux430ux43dux44fux435ux442ux441ux44f-ux438ux43dux444ux43eux440ux43cux430ux446ux438ux44f-ux43e-ux442ux43eux43c-ux43aux442ux43e-ux43aux43eux433ux434ux430-ux438-ux43aux430ux43aux438ux435-ux438ux43cux435ux43dux43dux43e-ux438ux437ux43cux435ux43dux435ux43dux438ux44f-ux432ux43dux435ux441-ux432-ux444ux430ux439ux43bux44b.-ux43fux440ux438-ux43dux435ux43eux431ux445ux43eux434ux438ux43cux43eux441ux442ux438-ux43cux43eux436ux43dux43e-ux43fux43eux441ux43cux43eux442ux440ux435ux442ux44c-ux43aux430ux43a-ux444ux430ux439ux43b-ux432ux44bux433ux43bux44fux434ux435ux43b-ux432-ux43fux440ux43eux448ux43bux43eux43c.}

\subsection{Возврат к предыдущим состояниям: Позволяет легко «откатить»
проект до любой точки в прошлом, если текущие изменения «сломали»
код.}\label{ux432ux43eux437ux432ux440ux430ux442-ux43a-ux43fux440ux435ux434ux44bux434ux443ux449ux438ux43c-ux441ux43eux441ux442ux43eux44fux43dux438ux44fux43c-ux43fux43eux437ux432ux43eux43bux44fux435ux442-ux43bux435ux433ux43aux43e-ux43eux442ux43aux430ux442ux438ux442ux44c-ux43fux440ux43eux435ux43aux442-ux434ux43e-ux43bux44eux431ux43eux439-ux442ux43eux447ux43aux438-ux432-ux43fux440ux43eux448ux43bux43eux43c-ux435ux441ux43bux438-ux442ux435ux43aux443ux449ux438ux435-ux438ux437ux43cux435ux43dux435ux43dux438ux44f-ux441ux43bux43eux43cux430ux43bux438-ux43aux43eux434.}

\subsection{Совместная работа: Дает возможность нескольким разработчикам
одновременно работать над одними и теми же файлами, а затем
автоматически объединять их
изменения.}\label{ux441ux43eux432ux43cux435ux441ux442ux43dux430ux44f-ux440ux430ux431ux43eux442ux430-ux434ux430ux435ux442-ux432ux43eux437ux43cux43eux436ux43dux43eux441ux442ux44c-ux43dux435ux441ux43aux43eux43bux44cux43aux438ux43c-ux440ux430ux437ux440ux430ux431ux43eux442ux447ux438ux43aux430ux43c-ux43eux434ux43dux43eux432ux440ux435ux43cux435ux43dux43dux43e-ux440ux430ux431ux43eux442ux430ux442ux44c-ux43dux430ux434-ux43eux434ux43dux438ux43cux438-ux438-ux442ux435ux43cux438-ux436ux435-ux444ux430ux439ux43bux430ux43cux438-ux430-ux437ux430ux442ux435ux43c-ux430ux432ux442ux43eux43cux430ux442ux438ux447ux435ux441ux43aux438-ux43eux431ux44aux435ux434ux438ux43dux44fux442ux44c-ux438ux445-ux438ux437ux43cux435ux43dux435ux43dux438ux44f.}

\subsection{Анализ изменений: Помогает понять причину появления ошибки
(бага), просматривая историю изменений конкретной строки
кода.}\label{ux430ux43dux430ux43bux438ux437-ux438ux437ux43cux435ux43dux435ux43dux438ux439-ux43fux43eux43cux43eux433ux430ux435ux442-ux43fux43eux43dux44fux442ux44c-ux43fux440ux438ux447ux438ux43dux443-ux43fux43eux44fux432ux43bux435ux43dux438ux44f-ux43eux448ux438ux431ux43aux438-ux431ux430ux433ux430-ux43fux440ux43eux441ux43cux430ux442ux440ux438ux432ux430ux44f-ux438ux441ux442ux43eux440ux438ux44e-ux438ux437ux43cux435ux43dux435ux43dux438ux439-ux43aux43eux43dux43aux440ux435ux442ux43dux43eux439-ux441ux442ux440ux43eux43aux438-ux43aux43eux434ux430.}

\subsection{Резервное копирование: Хранилище кода служит надежным
бэкапом
проекта.}\label{ux440ux435ux437ux435ux440ux432ux43dux43eux435-ux43aux43eux43fux438ux440ux43eux432ux430ux43dux438ux435-ux445ux440ux430ux43dux438ux43bux438ux449ux435-ux43aux43eux434ux430-ux441ux43bux443ux436ux438ux442-ux43dux430ux434ux435ux436ux43dux44bux43c-ux431ux44dux43aux430ux43fux43eux43c-ux43fux440ux43eux435ux43aux442ux430.}

\section{Объясните следующие понятия VCS и их отношения: хранилище,
commit, история, рабочая
копия.}\label{ux43eux431ux44aux44fux441ux43dux438ux442ux435-ux441ux43bux435ux434ux443ux44eux449ux438ux435-ux43fux43eux43dux44fux442ux438ux44f-vcs-ux438-ux438ux445-ux43eux442ux43dux43eux448ux435ux43dux438ux44f-ux445ux440ux430ux43dux438ux43bux438ux449ux435-commit-ux438ux441ux442ux43eux440ux438ux44f-ux440ux430ux431ux43eux447ux430ux44f-ux43aux43eux43fux438ux44f.}

Хранилище (Repository / Репозиторий): Это место, где VCS хранит все
файлы проекта и всю историю их изменений. Можно представить как
специальную базу данных или папку с проектом, за которой «следит» VCS.

Commit (Коммит / Фиксация): Это «снимок» (snapshot) состояния проекта на
определенный момент времени. Когда вы сделали какое-то осмысленное
изменение (например, добавили новую функцию или исправили ошибку), вы
создаете коммит. Он сохраняет текущее состояние файлов и сопровождается
комментарием от разработчика.

История (History / Лог): Это цепочка коммитов, расположенных в
хронологическом порядке. История позволяет проследить эволюцию проекта,
посмотреть, кто и что менял в каждом коммите.

Рабочая копия (Working Copy): Это текущее состояние файлов проекта на
вашем локальном компьютере, с которым вы непосредственно работаете.
Именно здесь вы редактируете код.

\section{Что представляют собой и чем отличаются централизованные и
децентрализованные VCS? Приведите
примеры.}\label{ux447ux442ux43e-ux43fux440ux435ux434ux441ux442ux430ux432ux43bux44fux44eux442-ux441ux43eux431ux43eux439-ux438-ux447ux435ux43c-ux43eux442ux43bux438ux447ux430ux44eux442ux441ux44f-ux446ux435ux43dux442ux440ux430ux43bux438ux437ux43eux432ux430ux43dux43dux44bux435-ux438-ux434ux435ux446ux435ux43dux442ux440ux430ux43bux438ux437ux43eux432ux430ux43dux43dux44bux435-vcs-ux43fux440ux438ux432ux435ux434ux438ux442ux435-ux43fux440ux438ux43cux435ux440ux44b.}

\subsection{Централизованные VCS
(CVCS):}\label{ux446ux435ux43dux442ux440ux430ux43bux438ux437ux43eux432ux430ux43dux43dux44bux435-vcs-cvcs}

\subsubsection{Принцип работы: Существует одно единственное главное
хранилище (сервер). Чтобы начать работу, разработчик копирует нужные
файлы из этого хранилища (делает чекаут), работает с ними, а затем
отправляет свои изменения обратно на сервер. Если сервер недоступен ---
вы не можете сохранить изменения, сравнить с историей или сделать
коммит.}\label{ux43fux440ux438ux43dux446ux438ux43f-ux440ux430ux431ux43eux442ux44b-ux441ux443ux449ux435ux441ux442ux432ux443ux435ux442-ux43eux434ux43dux43e-ux435ux434ux438ux43dux441ux442ux432ux435ux43dux43dux43eux435-ux433ux43bux430ux432ux43dux43eux435-ux445ux440ux430ux43dux438ux43bux438ux449ux435-ux441ux435ux440ux432ux435ux440.-ux447ux442ux43eux431ux44b-ux43dux430ux447ux430ux442ux44c-ux440ux430ux431ux43eux442ux443-ux440ux430ux437ux440ux430ux431ux43eux442ux447ux438ux43a-ux43aux43eux43fux438ux440ux443ux435ux442-ux43dux443ux436ux43dux44bux435-ux444ux430ux439ux43bux44b-ux438ux437-ux44dux442ux43eux433ux43e-ux445ux440ux430ux43dux438ux43bux438ux449ux430-ux434ux435ux43bux430ux435ux442-ux447ux435ux43aux430ux443ux442-ux440ux430ux431ux43eux442ux430ux435ux442-ux441-ux43dux438ux43cux438-ux430-ux437ux430ux442ux435ux43c-ux43eux442ux43fux440ux430ux432ux43bux44fux435ux442-ux441ux432ux43eux438-ux438ux437ux43cux435ux43dux435ux43dux438ux44f-ux43eux431ux440ux430ux442ux43dux43e-ux43dux430-ux441ux435ux440ux432ux435ux440.-ux435ux441ux43bux438-ux441ux435ux440ux432ux435ux440-ux43dux435ux434ux43eux441ux442ux443ux43fux435ux43d-ux432ux44b-ux43dux435-ux43cux43eux436ux435ux442ux435-ux441ux43eux445ux440ux430ux43dux438ux442ux44c-ux438ux437ux43cux435ux43dux435ux43dux438ux44f-ux441ux440ux430ux432ux43dux438ux442ux44c-ux441-ux438ux441ux442ux43eux440ux438ux435ux439-ux438ux43bux438-ux441ux434ux435ux43bux430ux442ux44c-ux43aux43eux43cux43cux438ux442.}

Примеры: CVS, Subversion (SVN).

\subsection{Децентрализованные / Распределенные VCS
(DVCS):}\label{ux434ux435ux446ux435ux43dux442ux440ux430ux43bux438ux437ux43eux432ux430ux43dux43dux44bux435-ux440ux430ux441ux43fux440ux435ux434ux435ux43bux435ux43dux43dux44bux435-vcs-dvcs}

\subsubsection{Принцип работы: Каждый разработчик клонирует (копирует)
себе не просто последнюю версию файлов, а полное хранилище (репозиторий)
целиком, вместе со всей историей. Таким образом, у каждого на компьютере
есть полноценная копия проекта. Большинство операций (просмотр истории,
создание коммитов) можно делать без доступа к сети. Сервер используется
для синхронизации и обмена изменениями между
разработчиками.}\label{ux43fux440ux438ux43dux446ux438ux43f-ux440ux430ux431ux43eux442ux44b-ux43aux430ux436ux434ux44bux439-ux440ux430ux437ux440ux430ux431ux43eux442ux447ux438ux43a-ux43aux43bux43eux43dux438ux440ux443ux435ux442-ux43aux43eux43fux438ux440ux443ux435ux442-ux441ux435ux431ux435-ux43dux435-ux43fux440ux43eux441ux442ux43e-ux43fux43eux441ux43bux435ux434ux43dux44eux44e-ux432ux435ux440ux441ux438ux44e-ux444ux430ux439ux43bux43eux432-ux430-ux43fux43eux43bux43dux43eux435-ux445ux440ux430ux43dux438ux43bux438ux449ux435-ux440ux435ux43fux43eux437ux438ux442ux43eux440ux438ux439-ux446ux435ux43bux438ux43aux43eux43c-ux432ux43cux435ux441ux442ux435-ux441ux43e-ux432ux441ux435ux439-ux438ux441ux442ux43eux440ux438ux435ux439.-ux442ux430ux43aux438ux43c-ux43eux431ux440ux430ux437ux43eux43c-ux443-ux43aux430ux436ux434ux43eux433ux43e-ux43dux430-ux43aux43eux43cux43fux44cux44eux442ux435ux440ux435-ux435ux441ux442ux44c-ux43fux43eux43bux43dux43eux446ux435ux43dux43dux430ux44f-ux43aux43eux43fux438ux44f-ux43fux440ux43eux435ux43aux442ux430.-ux431ux43eux43bux44cux448ux438ux43dux441ux442ux432ux43e-ux43eux43fux435ux440ux430ux446ux438ux439-ux43fux440ux43eux441ux43cux43eux442ux440-ux438ux441ux442ux43eux440ux438ux438-ux441ux43eux437ux434ux430ux43dux438ux435-ux43aux43eux43cux43cux438ux442ux43eux432-ux43cux43eux436ux43dux43e-ux434ux435ux43bux430ux442ux44c-ux431ux435ux437-ux434ux43eux441ux442ux443ux43fux430-ux43a-ux441ux435ux442ux438.-ux441ux435ux440ux432ux435ux440-ux438ux441ux43fux43eux43bux44cux437ux443ux435ux442ux441ux44f-ux434ux43bux44f-ux441ux438ux43dux445ux440ux43eux43dux438ux437ux430ux446ux438ux438-ux438-ux43eux431ux43cux435ux43dux430-ux438ux437ux43cux435ux43dux435ux43dux438ux44fux43cux438-ux43cux435ux436ux434ux443-ux440ux430ux437ux440ux430ux431ux43eux442ux447ux438ux43aux430ux43cux438.}

Примеры: Git, Mercurial.

\section{Опишите действия с VCS при единоличной работе с
хранилищем.}\label{ux43eux43fux438ux448ux438ux442ux435-ux434ux435ux439ux441ux442ux432ux438ux44f-ux441-vcs-ux43fux440ux438-ux435ux434ux438ux43dux43eux43bux438ux447ux43dux43eux439-ux440ux430ux431ux43eux442ux435-ux441-ux445ux440ux430ux43dux438ux43bux438ux449ux435ux43c.}

Инициализация: Создать новый проект и превратить его в репозиторий (git
init).

Работа: Создавать и редактировать файлы в рабочей копии.

Подготовка (Staging): Выбрать файлы, которые войдут в следующий коммит.
Это как собрать посылку --- вы кладете в нее только то, что нужно
отправить (git add).

Фиксация (Commit): «Сфотографировать» подготовленные файлы и сохранить
этот снимок в репозиторий, подписав его понятным сообщением (git commit
-m \enquote{Сообщение}).

Просмотр истории: Посмотреть список сделанных коммитов, чтобы вспомнить,
что и когда делалось (git log).

Возврат к прошлому: При необходимости откатить проект до состояния
одного из старых коммитов (git checkout или git revert).

\section{Опишите порядок работы с общим хранилищем
VCS.}\label{ux43eux43fux438ux448ux438ux442ux435-ux43fux43eux440ux44fux434ux43eux43a-ux440ux430ux431ux43eux442ux44b-ux441-ux43eux431ux449ux438ux43c-ux445ux440ux430ux43dux438ux43bux438ux449ux435ux43c-vcs.}

Получение актуальной версии: Перед началом работы разработчик забирает
себе последние изменения из общего репозитория (git pull).

Работа над своей задачей: Разработчик вносит изменения в своей локальной
рабочей копии, делая коммиты (git add, git commit).

Подготовка к отправке: Перед отправкой своих изменений в общий
репозиторий разработчик снова забирает последние изменения (git pull).
Это нужно на тот случай, пока он работал, коллеги уже что-то изменили.

Разрешение конфликтов: Если изменения коллег затрагивают те же строки
кода, что и ваши, может возникнуть конфликт. Разработчик вручную правит
файлы, оставляя нужный код, и фиксирует результат.

Отправка изменений: Теперь, когда локальная версия объединена с
последней версией из общего репозитория, можно отправить свои коммиты на
сервер (git push).

\section{Каковы основные задачи, решаемые инструментальным средством
git?}\label{ux43aux430ux43aux43eux432ux44b-ux43eux441ux43dux43eux432ux43dux44bux435-ux437ux430ux434ux430ux447ux438-ux440ux435ux448ux430ux435ux43cux44bux435-ux438ux43dux441ux442ux440ux443ux43cux435ux43dux442ux430ux43bux44cux43dux44bux43c-ux441ux440ux435ux434ux441ux442ux432ux43eux43c-git}

Управление историей разработки.

Обеспечение удобной работы с ветками.

Предоставление возможности каждому разработчику иметь полную копию
проекта.

Гарантия целостности кода.

Обеспечение нелинейной разработки.

\section{Назовите и дайте краткую характеристику командам
git.}\label{ux43dux430ux437ux43eux432ux438ux442ux435-ux438-ux434ux430ux439ux442ux435-ux43aux440ux430ux442ux43aux443ux44e-ux445ux430ux440ux430ux43aux442ux435ux440ux438ux441ux442ux438ux43aux443-ux43aux43eux43cux430ux43dux434ux430ux43c-git.}

git init --- Создать новый пустой репозиторий в текущей папке.

git clone --- Скопировать существующий удаленный репозиторий к себе на
компьютер.

git status --- Показать текущее состояние рабочей копии: какие файлы
изменены, какие готовы к коммиту.

git add --- Добавить файл в область подготовленных изменений (staging
area).

git commit -m \enquote{сообщение} --- Зафиксировать подготовленные
изменения в репозитории с комментарием.

git log --- Показать историю коммитов.

git push --- Отправить новые коммиты из локального репозитория в
удаленный.

git pull --- Забрать новые изменения из удаленного репозитория и
объединить их с локальной версией.

git branch --- Показать список веток или создать новую ветку.

git checkout --- Переключиться на другую ветку.

git merge --- Влить изменения из указанной ветки в текущую.

git diff --- Показать различия между файлами (например, между рабочей
копией и последним коммитом).

\section{Приведите примеры использования при работе с локальным и
удалённым
репозиториями.}\label{ux43fux440ux438ux432ux435ux434ux438ux442ux435-ux43fux440ux438ux43cux435ux440ux44b-ux438ux441ux43fux43eux43bux44cux437ux43eux432ux430ux43dux438ux44f-ux43fux440ux438-ux440ux430ux431ux43eux442ux435-ux441-ux43bux43eux43aux430ux43bux44cux43dux44bux43c-ux438-ux443ux434ux430ux43bux451ux43dux43dux44bux43c-ux440ux435ux43fux43eux437ux438ux442ux43eux440ux438ux44fux43cux438.}

Локальный :

mkdir my\_project cd my\_project git init echo
\enquote{print(\enquote*{Hello World})} \textgreater{} main.py git add
main.py git commit -m \enquote{Initial commit with main.py} git log
--oneline

Удалённый:

git clone https://github.com/username/project.git cd project git add
README.md git commit -m \enquote{Обновил инструкцию в README} git push
origin main git pull origin main

\section{Что такое и зачем могут быть нужны ветви
(branches)?}\label{ux447ux442ux43e-ux442ux430ux43aux43eux435-ux438-ux437ux430ux447ux435ux43c-ux43cux43eux433ux443ux442-ux431ux44bux442ux44c-ux43dux443ux436ux43dux44b-ux432ux435ux442ux432ux438-branches}

Ветка (branch) --- это отдельная линия разработки. По сути, это просто
перемещаемый указатель на один из коммитов. Создавая новую ветку, вы
создаете «копию» проекта на момент создания ветки и можете работать в
ней независимо.

Зачем нужны ветки:

Разработка новых функций (feature branches): Для каждой новой большой
функции создается отдельная ветка. Разработчик работает в ней, не боясь
сломать основной, стабильный код.

Исправление ошибок (bugfix branches): Срочное исправление бага можно
сделать в отдельной ветке, пока основная разработка новой функции
продолжается параллельно.

Эксперименты: Можно создать ветку для risky-эксперимента, и если он не
удастся --- просто удалить ветку, не повредив основной проект.

Поддержка разных версий: Можно вести разработку основной версии (main)
и, например, отдельно поддерживать ветку для версии 1.0, исправляя в ней
только критические баги.

\section{Как и зачем можно игнорировать некоторые файлы при
commit?}\label{ux43aux430ux43a-ux438-ux437ux430ux447ux435ux43c-ux43cux43eux436ux43dux43e-ux438ux433ux43dux43eux440ux438ux440ux43eux432ux430ux442ux44c-ux43dux435ux43aux43eux442ux43eux440ux44bux435-ux444ux430ux439ux43bux44b-ux43fux440ux438-commit}

Чтобы игнорировать файлы, нужно создать специальный файл с именем
.gitignore в корне репозитория. Внутри этого файла перечисляются шаблоны
имен файлов и папок, которые Git должен игнорировать.

Зачем это нужно:

Временные файлы: Файлы, создаваемые редактором или операционной системой
(например, .DS\_Store на macOS, *.swp в Vim).

Файлы с секретами: Файлы конфигурации, содержащие пароли, ключи API,
токены доступа.

Результаты сборки (бинарники): Скомпилированные файлы (.exe, .class,
.dll), которые можно заново собрать из исходного кода.

Зависимости и библиотеки: Папки вроде node\_modules/ (для JavaScript)
или vendor/ (для PHP). Они обычно очень большие и могут быть
восстановлены из файлов-манифестов (package.json, composer.json).

Логи и кэш: Файлы с логами работы программы или временные файлы кэша.

Хранить в репозитории нужно только исходный код и конфигурацию проекта,
чтобы он оставался легким, безопасным и чистым от мусора.

\chapter{Выводы}\label{ux432ux44bux432ux43eux434ux44b}

\begin{enumerate}
\def\labelenumi{\arabic{enumi})}
\tightlist
\item
  Изучили идеологию и применение средств контроля версий.
\item
  Освоили умения по работе с git.
\end{enumerate}

\chapter*{Список
литературы}\label{ux441ux43fux438ux441ux43eux43a-ux43bux438ux442ux435ux440ux430ux442ux443ux440ux44b}
\addcontentsline{toc}{chapter}{Список литературы}

\printbibliography[heading=none]





\end{document}
